\chapter*{Puente hacia la evidencia interna}
\addcontentsline{toc}{chapter}{Puente hacia la evidencia interna}

\section*{Separación de nomenclaturas}

\begin{observation}
La jerarquía dimensional D0--D11 pertenece al desarrollo canónico de MADELEINE. Las denominaciones \texttt{LAB-D0}, \texttt{LAB-D1} y \texttt{LAB-D2} pertenecen, en cambio, a una batería experimental interna posterior. El prefijo \texttt{LAB-} es obligatorio cuando se habla de esos controles y candidatos experimentales.
\end{observation}

\begin{observation}
No se establece equivalencia entre D0, D1 o D2 de la jerarquía dimensional y \texttt{LAB-D0}, \texttt{LAB-D1} o \texttt{LAB-D2}. La semejanza nominal identifica únicamente niveles de una secuencia de control dentro del laboratorio.
\end{observation}

\section*{Función de la evidencia posterior}

\begin{observation}
El fundamento presentado en este tomo no depende retrospectivamente del éxito de una batería numérica concreta. Los experimentos posteriores sirven para someter ciertas lecturas operativas a reglas explícitas de falsación, comparación, trazabilidad y estabilidad local.
\end{observation}

\begin{observation}
La exposición detallada de esa evidencia corresponde al Tomo III. Tomo I conserva su función fundacional y sólo fija aquí la separación entre canon, lectura operativa y evidencia interna.
\end{observation}

\section*{Cláusula de alcance}

\begin{observation}
La existencia de resultados internos reproducibles no convierte por sí sola a MADELEINE en una teoría física validada, no establece universalidad y no obliga a interpretar geométricamente los regímenes construidos.
\end{observation}
